% GENERATED FILE. Edit canonical lessons, then run npm run generate:books.
% canonical-source-hash: fnv1a64:4d7d785b599df7a0
% canonical-lessons: ZH-C339-changjia, ZH-C339-guoqing, ZH-W339-qing-celebrate, ZH-C339-qingzhu, ZH-C339-xiqing, ZH-C339-yanhua

\chapter{National Day}
\label{ch:zh-national-day}
\hlchaptermodality{339}

\begin{chapteropening}
\textbf{By the end of this chapter:} \emph{I can talk about a long holiday, National Day, celebrating, a festive house and fireworks.}

Complete the last lesson of chapter 339: I can talk about a long holiday, National Day, celebrating, a festive house and fireworks.
\end{chapteropening}

\section[\texorpdfstring{\zh{长假} (chángjià)}{长假 (chángjià)}]{\zh{长假} (chángjià) — long holiday}

\label{lesson:ZH-C339-changjia}



\begin{quote}
Before the new one: say the Mandarin for a day off work, then the Mandarin for a local custom.
\end{quote}

\subsection*{You'll want to know: \zh{长假}}
\textbf{\zh{长假}} — \emph{chángjià} — \textquotedblleft{}long holiday\textquotedblright{}.

\begin{quote}
\textbf{\zh{十月有一个长假}} \emph{shíyuè yǒu yí ge chángjià} — There's a long holiday in October.
\end{quote}

\begin{cousinweb}
\textbf{\zh{长假}} (\emph{chángjià}) is built from \zh{长} 'long' and \zh{假} 'leave', both inherited from Middle Chinese, \emph{drjang} and \emph{kaeH} (Guangyun, in Baxter's transcription), which glosses the first 'long' and 'far' and the second 'a holiday from work'.
\end{cousinweb}

\begin{grammarlens}[title={What you've built}]
One more word for telling what happened.
\end{grammarlens}

\subsection*{Your turn}
\begin{itemize}
\raggedright
  \item \emph{Say it:} \emph{chángjià}
  \item \emph{Say it:} \emph{chángjià}, once more
  \item \emph{Say it:} \emph{shíyuè yǒu yí ge chángjià} — There's a long holiday in October.
  \item \emph{Recall:} say the Mandarin for a day off work, then the Mandarin for a local custom, then say \emph{chángjià} again
\end{itemize}

\begin{tcolorbox}[breakable,colback=teal!4,colframe=teal!35!black,title={Before you move on}]
What does \zh{长假} mean? (\textquotedblleft{}Long holiday\textquotedblright{}.) Say it once more.
\end{tcolorbox}
\section[\texorpdfstring{\zh{国庆} (Guóqìng)}{国庆 (Guóqìng)}]{\zh{国庆} (Guóqìng) — National Day}

\label{lesson:ZH-C339-guoqing}



\begin{quote}
Before the new one: say the Mandarin for a local custom, then the Mandarin for a long holiday.
\end{quote}

\subsection*{You'll want to know: \zh{国庆}}
\textbf{\zh{国庆}} — \emph{Guóqìng} — \textquotedblleft{}National Day\textquotedblright{}.

\begin{quote}
\textbf{\zh{国庆的时候我们回了老家}} \emph{Guóqìng de shíhou wǒmen huí le lǎojiā} — Over the national holiday we went back to our home town.
\end{quote}

The new character, in pieces: \zh{庆} — six strokes: the roof piece of \zh{店}, a dot, a bar and a long stroke falling to the left; under it \zh{大}. The next lesson puts it in your hand.

\begin{cousinweb}
\textbf{\zh{国庆}} (\emph{Guóqìng}) is built from \zh{国} 'a state' and \emph{qìng} 'to celebrate', both inherited from Middle Chinese, \emph{kwok} and \emph{khjaengH} (Guangyun, in Baxter's transcription), which glosses the second 'to congratulate' and 'good fortune'.
\end{cousinweb}

\begin{grammarlens}[title={What you've built}]
One more word for telling what happened.
\end{grammarlens}

\subsection*{Your turn}
\begin{itemize}
\raggedright
  \item \emph{Say it:} \emph{Guóqìng}
  \item \emph{Say it:} \emph{Guóqìng}, once more
  \item \emph{Say it:} \emph{Guóqìng de shíhou wǒmen huí le lǎojiā} — Over the national holiday we went back to our home town.
  \item \emph{Recall:} say the Mandarin for a local custom, then the Mandarin for a long holiday, then say \emph{Guóqìng} again
\end{itemize}

\begin{tcolorbox}[breakable,colback=teal!4,colframe=teal!35!black,title={Before you move on}]
What does \zh{国庆} mean? (\textquotedblleft{}National Day\textquotedblright{}.) Say it once more.
\end{tcolorbox}
\section[\texorpdfstring{\zh{庆} (qìng)}{庆 (qìng)}]{\zh{庆} — the roof piece over \zh{大}}

\label{lesson:ZH-W339-qing-celebrate}



\begin{quote}
The word first, then the character inside it.

\begin{itemize}
\raggedright
  \item \emph{Recall:} say \emph{Guóqìng} — \textbf{R1}, one lesson back, before it decays
\end{itemize}

\begin{quote}
\zh{店} \zh{大}
\end{quote}

\emph{Read it:} each one aloud, with nothing said first — \emph{diàn}, \emph{dà}
\end{quote}

\subsection*{Script — the roof piece over \zh{大}}
\hlblockfigure{\detokenize{figures/ZH-W339-qing-celebrate-filmstrip.pdf}}{How \zh{庆} is written, stroke by stroke}

\begin{quote}
\zh{庆}
\end{quote}

\textbf{Two pieces: the roof first, then what sits under it.}

\textbf{The roof piece, three strokes,} as on \zh{店}: a dot, a bar, then a long stroke falling to the left down the side.

\textbf{The piece under it, \zh{大}, three strokes:} a bar, a stroke falling to the left through it, then a stroke falling to the right.

3 + 3 = six.

\subsection*{Your turn}
\begin{itemize}
\raggedright
  \item \emph{Write it:} \zh{庆}, all six, with the model in front of you
  \item \emph{Cover:} write it again from memory, and count the pieces, not the strokes
\end{itemize}

\begin{tcolorbox}[breakable,colback=teal!4,colframe=teal!35!black,title={Before you move on}]
How many strokes in \zh{庆}? (\textbf{Six}.) What is the roof? (\textbf{The roof piece of \zh{店}}.) What sits under it? (\textbf{\zh{大}}.)

Source: \href{https://raw.githubusercontent.com/chanind/hanzi-writer-data/68d10a4b21150cae5e1ebbd223eed289cf32d90c/data/\%E5\%BA\%86.json}{Hanzi Writer Data \zh{庆}, pinned stroke paths and medians}.
\end{tcolorbox}
\section[\texorpdfstring{\zh{庆祝} (qìngzhù)}{庆祝 (qìngzhù)}]{\zh{庆祝} (qìngzhù) — to celebrate}

\label{lesson:ZH-C339-qingzhu}



\begin{quote}
Before the new one: say the Mandarin for a long holiday, then the Mandarin for the National Day holiday.
\end{quote}

\subsection*{You'll want to know: \zh{庆祝}}
\textbf{\zh{庆祝}} — \emph{qìngzhù} — \textquotedblleft{}to celebrate\textquotedblright{}.

\begin{quote}
\textbf{\zh{我们一起庆祝国庆}} \emph{wǒmen yìqǐ qìngzhù Guóqìng} — We celebrate the national holiday together.
\end{quote}

\begin{cousinweb}
\textbf{\zh{庆祝}} (\emph{qìngzhù}) is built from \zh{庆} 'to congratulate' and \zh{祝} 'to bless', both inherited from Middle Chinese, \emph{khjaengH} and \emph{tsyuwH} (Guangyun, in Baxter's transcription), which glosses the first 'to congratulate'.
\end{cousinweb}

\begin{grammarlens}[title={What you've built}]
One more word for telling what happened.
\end{grammarlens}

\subsection*{Your turn}
\begin{itemize}
\raggedright
  \item \emph{Say it:} \emph{qìngzhù}
  \item \emph{Say it:} \emph{qìngzhù}, once more
  \item \emph{Say it:} \emph{wǒmen yìqǐ qìngzhù Guóqìng} — We celebrate the national holiday together.
  \item \emph{Recall:} say the Mandarin for a long holiday, then the Mandarin for the National Day holiday, then say \emph{qìngzhù} again
\end{itemize}

\begin{tcolorbox}[breakable,colback=teal!4,colframe=teal!35!black,title={Before you move on}]
What does \zh{庆祝} mean? (\textquotedblleft{}To celebrate\textquotedblright{}.) Say it once more.
\end{tcolorbox}
\section[\texorpdfstring{\zh{喜庆} (xǐqìng)}{喜庆 (xǐqìng)}]{\zh{喜庆} (xǐqìng) — festive; joyous}

\label{lesson:ZH-C339-xiqing}



\begin{quote}
Before the new one: say the Mandarin for the National Day holiday, then the Mandarin for to celebrate a happy event.
\end{quote}

\subsection*{You'll want to know: \zh{喜庆}}
\textbf{\zh{喜庆}} — \emph{xǐqìng} — \textquotedblleft{}festive; joyous\textquotedblright{}.

\begin{quote}
\textbf{\zh{过年的时候家里很喜庆}} \emph{guònián de shíhou jiā lǐ hěn xǐqìng} — At the new year the house is full of festive joy.
\end{quote}

\begin{cousinweb}
\textbf{\zh{喜庆}} (\emph{xǐqìng}) is built from \zh{喜} 'joy' and \zh{庆} 'good fortune', both inherited from Middle Chinese, \emph{xiX} and \emph{khjaengH} (Guangyun, in Baxter's transcription), which glosses the first 'joy' and the second 'good fortune'.
\end{cousinweb}

\begin{grammarlens}[title={What you've built}]
One more word for telling what happened.
\end{grammarlens}

\subsection*{Your turn}
\begin{itemize}
\raggedright
  \item \emph{Say it:} \emph{xǐqìng}
  \item \emph{Say it:} \emph{xǐqìng}, once more
  \item \emph{Say it:} \emph{guònián de shíhou jiā lǐ hěn xǐqìng} — At the new year the house is full of festive joy.
  \item \emph{Recall:} say the Mandarin for the National Day holiday, then the Mandarin for to celebrate a happy event, then say \emph{xǐqìng} again
\end{itemize}

\begin{tcolorbox}[breakable,colback=teal!4,colframe=teal!35!black,title={Before you move on}]
What does \zh{喜庆} mean? (\textquotedblleft{}Festive; joyous\textquotedblright{}.) Say it once more.
\end{tcolorbox}
\section[yānhuā]{yānhuā — fireworks}

\label{lesson:ZH-C339-yanhua}



\begin{quote}
Before the new one: say the Mandarin for to celebrate a happy event, then the Mandarin for full of festive joy.
\end{quote}

\subsection*{You'll want to know: yānhuā}
\emph{yānhuā} — \textquotedblleft{}fireworks\textquotedblright{}.

\begin{quote}
\emph{Guóqìng de wǎnshang wǒmen kàn le yānhuā} — On the evening of the national holiday we saw the fireworks.
\end{quote}

Its characters come later; for now it is a word for the ear.

\begin{cousinweb}
\emph{yānhuā} is built from \emph{yān} 'smoke' and \zh{花} 'a flower', both inherited from Middle Chinese, \emph{'en} and \emph{xwae} (Guangyun, in Baxter's transcription), which glosses the first 'the breath of fire'.
\end{cousinweb}

\begin{grammarlens}[title={What you've built}]
One more word for telling what happened.
\end{grammarlens}

\subsection*{Your turn}
\begin{itemize}
\raggedright
  \item \emph{Say it:} \emph{yānhuā}
  \item \emph{Say it:} \emph{yānhuā}, once more
  \item \emph{Say it:} \emph{Guóqìng de wǎnshang wǒmen kàn le yānhuā} — On the evening of the national holiday we saw the fireworks.
  \item \emph{Recall:} say the Mandarin for to celebrate a happy event, then the Mandarin for full of festive joy, then say \emph{yānhuā} again
\end{itemize}

\begin{tcolorbox}[breakable,colback=teal!4,colframe=teal!35!black,title={Before you move on}]
What does yānhuā mean? (\textquotedblleft{}Fireworks\textquotedblright{}.) Say it once more.
\end{tcolorbox}
